\chapter{Conclusion}

\section{Travaux réalisés}
Ce travail a permis de reprendre le démonstrateur là où le précédent travail de bachelor l'avait laissé, puis de le faire progresser sur plusieurs plans.

Dans un premier temps, le fonctionnement du matériel a été vérifié : circuit de précharge, convertisseurs DC/DC, capteurs de courant et de position. Ces mesures ont permis de valider l'électronique de la carte principale et de corriger un défaut de soudure sur le convertisseur 24~V.

La masse réelle de la partie lévitante a ensuite été mesurée (18.052~kg au lieu des 13.6~kg supposés), et la valeur des inductances a été redéterminée au RLC-mètre, révélant un écart important avec la valeur théorique. Ces deux corrections ont directement alimenté la synthèse des régulations.

Les circuits de conditionnement ont été modifiés pour exploiter toute la plage de l'ADC, et un filtre RC a été ajouté sur les mesures de position. La plaque de référence, découverte ondulée, a été remplacée par une plaque en composite d'aluminium bien plus régulière.

Côté logiciel, la régulation de courant a été reprise et une erreur de synthèse présente depuis plusieurs travaux a été corrigée. Surtout, une régulation de position en MIMO a été mise en place, en remplacement de l'approche SISO : reconstruction modale de la position, retour d'état avec observateur, allocation des efforts sur les quatre inducteurs et anti-windup. Une nouvelle machine d'état gère l'ensemble de la séquence de décollage, et l'interface web embarquée a été complétée avec l'affichage du courant. Enfin, la maquette a été rendue autonome grâce à l'ajout d'une alimentation propre.

\section{Retour sur les hypothèses}
Les hypothèses de recherche formulées dans la \autoref{sec:Hypothese} sont reprises ci-dessous, chacune accompagnée de son verdict à l'issue du projet :

\begin{itemize}
    \item Les interruptions liées à la communication sans fil perturbent l'interruption de régulation. Invalidée : les mesures d'exécution du programme montrent que l'interruption tient largement dans la période de 40~$\mu$s, avec ou sans la communication UART.

    \item Les contraintes temporelles ne sont plus respectées. Invalidée pour la même raison : la durée d'exécution des tâches laisse une marge confortable.

    \item L'ajout d'une action intégrale sur l'inducteur 3 engendre une instabilité. Partiellement validée : un comportement anormal a bien été observé sur cet inducteur, sans qu'il soit possible de le désigner comme cause unique du problème.

    \item Une mauvaise calibration globale empêche le décollage. Validée : la plaque de référence, découverte ondulée, faussait les mesures de position. Son remplacement a nettement amélioré la situation.

    \item L'intégration de l'électronique embarquée a modifié la masse totale. Validée : la pesée donne 18.052~kg au lieu des 13.6~kg supposés, ce qui fausse les paramètres de la régulation.

    \item Une erreur a été introduite lors de la synthèse des régulateurs. Validée : la constante $K_{ui}$ avait été omise dans la synthèse de la régulation de courant (\autoref{subsec:RegPIError}).

    \item La condition de transition à 99~\% de la machine d'état n'est pas robuste. Non testée : cette piste n'a pas été approfondie, la machine d'état ayant par ailleurs été entièrement revue.

    \item La carte principale présente un défaut matériel. Partiellement validée : une broche mal soudée sur le convertisseur 24~V a été trouvée et corrigée, mais le reste de l'électronique s'est révélé sain.

    \item Les broches du DSP ne sont pas correctement reliées à la carte principale. Non testée : cette hypothèse n'a pas été vérifiée, faute de temps.

    \item L'inductance nominale calculée est incorrecte. Validée : les mesures au RLC-mètre donnent entre 7.7~mH et 8.65~mH contre 4.9~mH en théorie.

    \item L'erreur n'est pas calculée correctement pour la boucle de position. Partiellement validée : la mesure du programme en fonctionnement a révélé une erreur transmise incohérente sur l'inducteur 1, sans que cela explique à lui seul le problème de lévitation.
\end{itemize}

\section{Améliorations futures}
La principale amélioration à apporter à la maquette serait de refaire la synthèse de la régulation de position MIMO. Celle-ci devrait être synthétisée par un véritable régulateur LQI, au moyen de la méthode de Riccati \cite{Riccati}, plutôt que par le placement de pôles empirique actuel.

Il faudrait ensuite trouver un moyen d'exécuter correctement le code depuis la flash, sans les problèmes d'oscillations rencontrés.

Enfin, l'interface utilisateur pourrait être enrichie. Elle ne sert actuellement qu'à afficher la position et le courant, et à enclencher ou couper la maquette. Il serait intéressant d'y ajouter des fonctions permettant d'interagir davantage avec la maquette.

L'ordre de priorité serait dès lors le suivant :

\begin{itemize}
    \item Réaliser une meilleure synthèse de la régulation de position.
    \item Améliorer l'exécution du code depuis la flash.
    \item Ajouter des fonctions d'interaction sur l'interface utilisateur.
\end{itemize}

\section{Conclusion}
Ce travail a été intellectuellement très enrichissant. Il m'a permis de consolider mes bases en mécatronique et d'aborder des notions que je n'avais jamais mises en pratique jusqu'ici.

Il m'a également permis de mettre à l'épreuve les différentes compétences acquises durant mes trois années de bachelor.

Bien que le résultat ne soit pas encore pleinement satisfaisant, ce travail constitue une avancée importante dans le projet.

\vfil
\hspace{8cm}\makeatletter\@author\makeatother\par
\hspace{8cm}\begin{minipage}{5cm}
    % Place pour signature numérique
    \printsignature
\end{minipage}